\documentclass[12pt]{article}
\usepackage[utf8]{inputenc}
\usepackage[T1]{fontenc}
\usepackage{hyperref}
\title{Upper Umpqua Language Resources for K-12 Curriculum Development}
\author{Compiled by Claude Sonnet 5.5 (claude-sonnet-5-5) for Andr\'as Kornai}
\date{}
\begin{document}
\maketitle

\section{Introduction}
Upper Umpqua (ISO 639-3 xup; Glottocode uppe1436) is an extinct Pacific Coast Athabaskan language of southwestern Oregon. Search results indicate that the record is fragmentary. It consists of Horatio Hale's 1841 vocabulary and notes collected from the last speakers by Melville Jacobs and John P. Harrington as late as the 1940s. I found no language-revitalization programme, school curriculum, or learner materials specific to Upper Umpqua. The resources below are background only.

\section{Primary Educational Resources}
No primers, textbooks or school materials for Upper Umpqua were found. A K-12 programme would have to be built from the archival sources listed under Linguistic and Reference Materials.

\section{Linguistic and Reference Materials}
\begin{itemize}
\item Wikipedia, ``Upper Umpqua language'': an overview that names the main sources (Hale 1841; Jacobs and Harrington fieldwork). \url{https://en.wikipedia.org/wiki/Upper_Umpqua_language}
\item Endangered Languages Project entry (the search result was the Siuslaw/Lower Umpqua entry, a different language, so it is only context). \url{https://endangeredlanguages.com/node/40353}
\end{itemize}

\section{Religious and Community Resources}
\begin{itemize}
\item Wikipedia, ``Umpqua people'': background on the Umpqua peoples and the Cow Creek Band (federally recognized 1982). The results did not show any current Cow Creek language programme. \url{https://en.wikipedia.org/wiki/Umpqua_people}
\item EBSCO research starter on the Umpqua: general background. \url{https://ebsco.com/research-starters/social-sciences-and-humanities/umpqua}
\end{itemize}

\section{Digital and Audio Resources}
No audio or digital resources for Upper Umpqua were found in search results. The Siuslaw/Lower Umpqua language is distinct and is not Upper Umpqua; its documentation should not be substituted for it.

\section{Curriculum Development Context}
Little exists for Upper Umpqua. Any curriculum would be reconstruction-based and would need to start from the archival manuscripts (Hale, Jacobs, Harrington) in consultation with the Cow Creek Band of Umpqua Tribe of Indians and neighbouring Oregon Athabaskan communities. I did not verify that any of those groups is interested in this work.

\end{document}